\input{preamble}

\section*{Dirac spinor normalization}

A widely-used convention is
\begin{equation*}
u^\dag u=2E
\end{equation*}

where $E$ is relativistic energy
\begin{equation*}
E=\sqrt{\vec p\cdot\vec p+m^2}
\end{equation*}

For example, for a spin-up electron
\begin{equation*}
u=\sqrt{E+m}\begin{pmatrix}
1
\\[1ex]
0
\\[1ex]
{\displaystyle\frac{p_z}{E+m}}
\\[2ex]
{\displaystyle\frac{p_x+ip_y}{E+m}}
\end{pmatrix}
\qquad\qquad
u^\dag u=2E
\tag{1}
\end{equation*}

To demonstrate Lorentz invariance, let $\vec p=0$.
Then $\sqrt{E+m}=\sqrt{2m}$ and
\begin{equation*}
u=\sqrt{2m}\begin{pmatrix}
1\\0\\0\\0
\end{pmatrix}
\qquad\qquad
u^\dag u=2m=2E
\end{equation*}

It is instructive to see what happens when $u$ is not normalized.
\begin{equation*}
u=\begin{pmatrix}
1
\\[1ex]
0
\\[1ex]
{\displaystyle\frac{p_z}{E+m}}
\\[2ex]
{\displaystyle\frac{p_x+ip_y}{E+m}}
\end{pmatrix}
\qquad\qquad
u^\dag u=\frac{\vec p\cdot\vec p}{(E+m)^2}+1
\tag{2}
\end{equation*}

The issue in (2) is $u^\dag u$ is not Lorentz invariant.
For example, for $\vec p=0$ we have
\begin{equation*}
u^\dag u=1
\end{equation*}

\href{https://georgeweigt.github.io/examples/dirac-spinor-normalization-demo.html}{Eigenmath script}

\end{document}
